% !TeX spellcheck = it_IT

\documentclass[handout]{beamer}

\usepackage[utf8]{inputenc}
\usepackage[italian]{babel}
\usepackage{multicol}
\usetheme{metropolis}

\usepackage{listings}

\usepackage{tikz}
\usetikzlibrary{shapes,arrows,positioning}

\tikzstyle{startstop} = [ellipse, draw, fill=bg, text=structure, text width=3em, text badly centered, inner sep=0pt]
\tikzstyle{block} = [rectangle, draw, fill=bg, text=structure, text width=5em, text centered]
\tikzstyle{decision} = [diamond, aspect=2, draw, fill=bg, text=structure, text width=3em, text badly centered, inner sep=0pt]
\tikzstyle{input} = [rectangle, draw, fill=bg, text=structure, text width=4em, text centered, rounded rectangle, rounded rectangle east arc=none]
\tikzstyle{output} = [rectangle, draw, fill=bg, text=structure, text width=4em, text centered, rounded rectangle, rounded rectangle west arc=none]

\tikzstyle{line} = [draw, -latex']


\title{Informatica B - Esercitazione 6}
\subtitle{Matlab}

\author{Stefano Cereda\\
		stefano1.cereda@mail.polimi.it
	}
\date{7/11/2017}
\institute[PoliMi]{\vspace{0.5cm}\centering Politecnico di Milano \\ \vspace{0.2cm}
	\includegraphics[width=0.2\linewidth]{./logopolimi}}
%\logo{\includegraphics[width=15mm]{./logopolimi}}

\setbeamercovered{invisible}

\makeindex

\begin{document}
\begin{frame}
	\maketitle
\end{frame}

\begin{frame}{Esercizio 0}
Date le seguenti matrici $a$ e $b$ si dica cosa stampano le seguenti istruzioni.
$a$ =
\begin{tabular}{|ccc|}
	\hline
	1&2&3\\
	4&5&6\\
	7&8&9\\
	\hline
\end{tabular}
$b$ =
\begin{tabular}{|ccc|}
	\hline
	9&8&7\\
	6&5&4\\
	3&2&1\\
	\hline
\end{tabular}

\begin{enumerate}
	\item a(3,1)
	\item a(8)
	\item mod(a,b)
	\item {[r c]}=size(mod(a,b))
\end{enumerate}

\end{frame}

\begin{frame}{Esercizio 1 - Rotazione di matrice}
Sc scriva una script in Matlab/Octave che data una matrice $a$ di dimensioni $N\times N$ crea una nuova matrice $b$ ruotata di 90 gradi in senso antiorario rispetto ad $a$.

Si consideri $N=4$ e la matrice $a$ inizializzata con i valori:

\centering
\begin{tabular}{|cccc|}
	\hline
	1&2&3&4\\
	2&3&4&5\\
	6&7&8&9\\
	0&0&0&0\\
	\hline
\end{tabular}
\end{frame}

\begin{frame}{Esercizio 2 - Sequenza palindroma}
Si scriva uno script che chiede l’inserimento di una sequenza di N numeri e alla fine dice se la sequenza è palindroma o no.

Si ricorda che una sequenza è palindroma se risulta uguale letta da sinistra verso destra o da destra verso sinistra
\end{frame}

\begin{frame}{Esercizio 3 - Confronti}
Si scriva uno script che chiede di inserire una sequenza di N numeri.

Successivamente, si inserisca un ulteriore numero e si dica se tutti i numeri della sequenza sono minori, uguali o maggiori di tale numero.
\end{frame}

\begin{frame}{Esercizio 4 - Confronti con Matrici}
Scrivere in Matlab uno script che, data una matrice $m$ di numeri, restituisce in uscita una matrice $mr$, ottenuta da $m$ nel seguente modo:
\begin{itemize}
	\item si calcola la media aritmetica dei valori di $m$;
	\item per i valori che in $m$ sono minori della media, in $m$r si
pone nella stessa posizione il valore -1, per quelli superiori alla media si pone il valore 1, e per gli altri (quelli uguali alla media) si pone lo stesso valore.
\end{itemize}

Si utilizzi la seguente matrice $m$

\centering
\begin{tabular}{|ccc|}
	\hline
	-2&2&3\\
	4&5&6\\
	\hline
\end{tabular}
\end{frame}

\begin{frame}{Esercizio 5 - Inserimento ordinato nell'array}
Scrivere un programma che acquisisce da tastiera una serie di interi e li memorizza in modo ordinato nell’array.

\pause
Se, per esempio, $a = \{1, 3, 4, …\}$ e si vuole inserire 2, i valori 3 e 4 devono essere spostati a destra per fare spazio al nuovo elemento:

\centering
\begin{tabular}{|c|c|c|c|}
	\hline
	1&3&4&...\\
	\hline
\end{tabular}
$\rightarrow$
\begin{tabular}{|c|c|c|c|c|}
	\hline
	1& \alert{2} &3&4&...\\
	\hline
\end{tabular}

\end{frame}

\end{document}